% Fichier généré automatiquement par tools/generate_corrections.py.
% Source : DYN/DYN-06-PFD/03_TT/03_TT.tex
% Statut : complet (2/2 question(s) renseignée(s), 0 reprise(s)).
\begin{corrigebox}[Corrigé extrait de la banque]
\CorrigeQuestion{1}
On isole \textbf{2}. 

\vspace{.25cm}

Bilan des actions mécaniques : 
\begin{itemize}
\item liaison glissière entre 1 et 2 : 
$\torseurstat{T}{1}{2}= \torseurl{X_{12}\vi{0}+Z_{12}\vk{0}}{L_{12}\vi{0}+M_{12}\vj{0}+N_{12}\vk{0}}{B}$;
\item pesanteur : $\torseurstat{T}{\text{Pes}}{2}= \torseurl{-m_2 g \vj{0}}{\vect{0}}{C}$;
\item vérin : $\torseurstat{T}{1_v}{2}= \torseurl{F_2 \vj{0}}{\vect{0}}{B}$.
\end{itemize}

\vspace{.25cm}

Application du TRD au solide \textbf{2} en projection sur $\vj{0}$: 

$\vectf{1}{2}\cdot \vj{0} + \vectf{\text{Pes}}{2}\cdot \vj{0}+ \vectf{1_v}{2}\cdot \vj{0} = \vectrd{2}{0}\cdot \vj{0}$.

\vspace{.25cm}

Calcul de la résultante dynamique : 
$\vectrd{2}{0} = m_2 \vectg{C}{2}{0} = m_2 \left(\lambdapp (t) \vi{0}+\mupp (t) \vj{0}\right)$.

Application du théorème : 
$$
-m_2g  + F_2 = m_2\mupp (t).
$$
\CorrigeQuestion{2}
On isole \textbf{1+2}. 

\vspace{.25cm}

Bilan des actions mécaniques : 
\begin{itemize}
\item liaison glissière entre 0 et 1 : 
$\torseurstat{T}{0}{1}= \torseurl{Y_{01}\vj{0}+Z_{12}\vk{0}}{L_{12}\vi{0}+M_{12}\vj{0}+N_{12}\vk{0}}{A}$;
\item pesanteur : $\torseurstat{T}{\text{Pes}}{1}= \torseurl{-m_1 g \vj{0}}{\vect{0}}{B}$;
\item pesanteur : $\torseurstat{T}{\text{Pes}}{2}= \torseurl{-m_2 g \vj{0}}{\vect{0}}{C}$;
\item vérin : $\torseurstat{T}{0_v}{1}= \torseurl{F_1 \vi{0}}{\vect{0}}{B}$.
\end{itemize}

\vspace{.25cm}

Application du TRD au solide \textbf{1+2} en projection sur $\vi{0}$: 

$\vectf{0}{1}\cdot \vi{0} + \vectf{\text{Pes}}{1}\cdot \vi{0}+ \vectf{\text{Pes}}{2}\cdot \vj{0}+ \vectf{0_v}{2}\cdot \vi{0} = \vectrd{1+2}{0}\cdot \vi{0}$.

\vspace{.25cm}

Calcul de la résultante dynamique : 
$\vectrd{1+2}{0} = m_1 \vectg{B}{1}{0}+m_2 \vectg{C}{2}{0} = m_1 \lambdapp (t) \vi{0} + m_2 \left(\lambdapp (t) \vi{0}+\mupp (t) \vj{0}\right)$.

Application du théorème : 
$$
F_1  + F_2 = m_1 \lambdapp (t)  + m_2 \lambdapp (t) .
$$
\end{corrigebox}
